\documentclass[11pt]{article}
\usepackage[margin=1in]{geometry}
\usepackage{amsmath,amssymb}
\usepackage{bussproofs}
\usepackage{enumitem}

\newcommand{\unit}{\mathsf{unit}}
\newcommand{\bool}{\mathsf{bool}}
\newcommand{\true}{\mathsf{true}}
\newcommand{\false}{\mathsf{false}}
\newcommand{\V}{\mathcal{V}}
\newcommand{\E}{\mathcal{E}}
\newcommand{\G}{\mathcal{G}}
\newcommand{\ok}{\;\mathsf{ok}}

\title{Homework 9 - Alex Bastianini}
\date{
  \today
}

\begin{document}
\maketitle

%--------------------------------------------------
\section*{1)}
\begin{enumerate}[label=\alph*.]
  \item
  \begin{itemize}
    \item $R_{\unit \to \unit \to \unit}(t_1)$ holds
    \item $t_2$ \texttt{() ()} is of type $\unit$ but fails to terminate
    \item $t_3$ \texttt{()} is of type $\unit \to \unit$ but fails to terminate
    \item $t_4$ is of type $\unit \to \unit \to \unit$ but fails to terminate
    \item $R(t_5)$ holds
  \end{itemize}

  \item $t_1, t_2, t_5$. $R_\tau$ is defined by induction on the type, so requiring only ``\ldots'' would include $t_2$ in the relation, when in reality $t_2$ \texttt{() ()} doesn't terminate.
\end{enumerate}

%--------------------------------------------------
\section*{2)}
\begin{enumerate}[label=\alph*.]
  \item
  \begin{align*}
    &\lambda f : \bool \to \bool.\\
    &\quad \textbf{if } f\ \true \textbf{ then}\\
    &\qquad \textbf{if } f\ \true \textbf{ then } \true\\
    &\qquad \textbf{else } \bot\\
    &\quad \textbf{else}\\
    &\qquad \textbf{if } f\ \true \textbf{ then } \bot\\
    &\qquad \textbf{else } \true
  \end{align*}

  \item $g \triangleq \lambda b : \bool.\ \false$

  \[
    \emptyset \not\vDash e_2 : (\bool \to \bool) \to \bool
    \iff
    \exists \gamma \in \G[\![\emptyset]\!].\ \gamma(e_2) \notin \E[\![(\bool \to \bool) \to \bool]\!]
  \]

  Let $\gamma = \emptyset$ (only possible).
  We need to prove that $e_2 \notin \V[\![\ldots]\!]$.
  Because $g \in \V[\![\bool \to \bool]\!]$, it suffices to prove that
  $e_0\{g/f\} \notin \E[\![\bool]\!]$
  (where $e_0$ is the body of $e_2$).

  We have that $e_0\{g/f\} \to^{*} \texttt{()()} \notin \V[\![\bool]\!]$.

  \item
  \[
    e_3 \triangleq \lambda f : \bool \to \bool.\ \textbf{if } (f\ \true \ \mathsf{OR}\ \true)\ \ldots
  \]
\end{enumerate}

%--------------------------------------------------
\section*{3)}
\begin{enumerate}[label=\arabic*.]
  \item None

  \item The two Church Booleans.

  \begin{prooftree}
    \AxiomC{$X \in \{X\}$}
    \UnaryInfC{$\{X\} \vdash X \ok$}
    \AxiomC{$\Gamma(a) = X$}
    \UnaryInfC{$\{X\}, \{a:X, b:X\} \vdash a : X$}
    \BinaryInfC{$\{X\}, \{a:X\} \vdash \lambda b{:}X.\ a : X \to X$}
    \UnaryInfC{$\{X\}, \emptyset \vdash \lambda a{:}X.\ \lambda b{:}X.\ a : X \to X \to X$}
    \UnaryInfC{$\emptyset, \emptyset \vdash \Lambda X.\ \lambda a{:}X.\ \lambda b{:}X.\ a : \forall X.\ X \to X \to X$}
  \end{prooftree}

  (similar for $\false$)

  We can't have any other terms since, as proven during lectures, a function with this type signature can only always return the first or only return the second argument.

  \item $\Lambda X.\ \Lambda Y.\ \lambda x{:}X.\ \lambda y{:}Y.\ x$

  \begin{prooftree}
    \AxiomC{$\{X,Y\} \vdash X, Y \ok$}
    \AxiomC{$\Gamma(x) = X$}
    \UnaryInfC{$\{X,Y\}, \{x:X, y:Y\} \vdash x : X$}
    \BinaryInfC{$\{X,Y\}, \emptyset \vdash \lambda x{:}X.\ \lambda y{:}Y.\ x : X \to Y \to X$}
    \UnaryInfC{$\vdash \Lambda X.\ \Lambda Y.\ \lambda x{:}X.\ \lambda y{:}Y.\ x : \forall X\, \forall Y.\ X \to Y \to X$}
  \end{prooftree}

  This is the only term, for the same reason for which $\forall X.\ X \to X$ is the identity.

  \item
  \begin{align*}
    &\Lambda X.\ \lambda x{:}X.\ \true\\
    &\Lambda X.\ \lambda x{:}X.\ \false
  \end{align*}

  \begin{prooftree}
    \AxiomC{$\{X\} \vdash X \ok$}
    \AxiomC{$\{X\}, \{x:X\} \vdash \true : \bool$}
    \BinaryInfC{$\{X\}, \emptyset \vdash \lambda x{:}X.\ \true : X \to \bool$}
    \UnaryInfC{$\vdash \Lambda X.\ \lambda x{:}X.\ \true : \forall X.\ X \to \bool$}
  \end{prooftree}

  We can't access the value of the argument, so we can either return $\true$ or $\false$.

  \item[5.] \phantom{x}
  \begin{enumerate}[label=\alph*.]
    \item
    \begin{align*}
      &\Lambda X.\ \lambda f{:}X \to X.\ \lambda x{:}X.\ x\\
      &\phantom{\Lambda X.\ \lambda f{:}X \to X.\ \lambda x{:}X.\ } f\ x\\
      &\phantom{\Lambda X.\ \lambda f{:}X \to X.\ \lambda x{:}X.\ } f\ (f\ x)\\
      &\phantom{\Lambda X.\ \lambda f{:}X \to X.\ \lambda x{:}X.\ } \vdots
    \end{align*}
    There are infinitely many terms, corresponding to the Church numerals.

    \item
    \begin{itemize}
      \item $R_1$: (i), (ii), (iii) hold, it works.
      \item $R_2$: none hold: $(0,0)$, $(2,7)$, $(1,4)$; doesn't work.
      \item $R_3$: (i), (iii) hold: $(11,33)$; doesn't work.
    \end{itemize}

    \item $R = \{(5n, 5n) \mid n \in \mathbb{Z}\}$

    By A.T.\ and (i), (ii), the results of two identical runs of the expression are related, so the expression evaluates to $5n$ for some $n \in \mathbb{Z}$, so it is a multiple of 5.
  \end{enumerate}
\end{enumerate}

%--------------------------------------------------
\section*{4)}
\begin{enumerate}[label=\alph*.]
  \item $R = \{(a, -a) \mid a \in \mathbb{Z}\}$

  \item
  \begin{align*}
    \mathit{client} \triangleq{}& \textbf{unpack } \{C, c\} = \mathit{counter} \textbf{ in}\\
    & (\#3\ c)\ \big((\#2\ c)\ ((\#2\ c)\ (\#1\ c))\big)
  \end{align*}

  \begin{itemize}
    \item In implementation $a$, $\mathit{client} \to^{*} 2$
    \item In implementation $c$, $\mathit{client} \to^{*} 0$
  \end{itemize}

  So $a$ and $c$ are not indistinguishable.
\end{enumerate}

\subsubsection*{Debriefing}
AI was used only for translating written solutions to LaTeX.

\begin{enumerate}
\item 4h
\item moderate
\item yes
\item 85
\end{enumerate}

\end{document}
